\section{预备知识}
\label{sec:2}

本节提供效应与协效应系统的简要概述——这是支撑我们工作的两个理论支柱。我们假设读者熟悉基本的类型论和范畴论；此处的目标是确定记号并引入第 \S\ref{sec:3} 节将作为运行时机制加以实现的关键抽象。

\subsection{效应}
\label{sec:2.1}

在简单类型 $\lambda$ 演算 (STLC) \cite{ref20, ref21} 中，类型推导 $\Gamma \vdash t : T$ 表明项 $t$ 在上下文 $\Gamma$ 下具有类型 $T$。效应系统对类型加以精化，以描述计算可能产生哪些副作用，从而得到如下形式的推导
\begin{equation}
\label{eq:2.1}
\Gamma \vdash t : T \mid \textit{effect}
\end{equation}
此处，结果类型附带效应代数的一个元素作为标注，该元素描述了计算可能产生哪些副作用，从而使得对有状态计算的组合推理成为可能。这一方法源自 Lucassen 和 Gifford \cite{ref22}，他们引入了一种区分类型、效应和区域的类别化类型系统，以发现并行程序中的调度约束。

\subsubsection{单子效应}
\label{sec:2.1.1}

Moggi \cite{ref16} 首次通过单子对计算效应进行范畴论建模；Wadler \cite{ref23} 在 Haskell 中推广了该方法。范畴 $\mathcal{C}$ 上的单子 $(T, \eta, \mu)$ 将效应计算封装为 $T(A)$ 类型的值，其中 $\eta : A \to T(A)$ 将纯值提升，$\mu : T(T(A)) \to T(A)$ 将嵌套计算顺序组合。经典实例包括 Maybe 单子（用于部分性）、State 单子（用于可变状态）和 IO 单子（用于外部交互）。

\subsubsection{代数效应}
\label{sec:2.1.2}

Plotkin 和 Power \cite{ref17, ref24} 证明了代数运算确定单子，建立了一个效应接口与其实现解耦的框架。效应签名 $\Sigma$ 声明一组运算（例如，用于状态的 $\mathrm{get} : () \to S$、$\mathrm{put} : S \to ()$）；程序自由地调用运算，而无需绑定到特定的解释。Plotkin 和 Pretnar \cite{ref25} 随后引入了效应处理器，通过提供续延语义来解释运算：
\begin{equation}
\label{eq:2.2}
\mathrm{handle}\ e\ \mathrm{with}\ \{ \mathrm{op}(v, \kappa) \mapsto \ldots \}
\end{equation}
处理器接收运算参数 $v$ 和定界续延 $\kappa$，它可以调用零次、一次或多次，从而在统一框架内实现异常、协程和非确定性 \cite{ref26}。诸如 Koka \cite{ref27, ref28}、Eff \cite{ref29} 和 OCaml 5 \cite{ref30} 等语言已在不同设计权衡下采用了代数效应。

\subsection{协效应}
\label{sec:2.2}

与效应对偶地，协效应系统 \cite{ref18, ref31} 丰富的是上下文而非类型，从而得到如下形式的推导
\begin{equation}
\label{eq:2.3}
\Gamma \mid \textit{coeffect} \vdash t : T
\end{equation}
此处，上下文附带协效应代数的一个元素作为标注，描述计算从其环境中需要什么，例如要访问的资源、要持有的权限或要依赖的服务。效应建模的是程序对世界的影响，而协效应建模的是世界对程序的约束。

\subsubsection{余单子协效应}
\label{sec:2.2.1}

使用余单子来构造上下文依赖计算的思想首先由 Uustalu 和 Vene \cite{ref32} 发展，他们提出了对称（半）单子余单子作为 Moggi 效应单子框架的对偶，捕获了数据流和属性求值等概念。Petricek 等人 \cite{ref18} 在此基础上提出了协效应，作为上下文依赖性的统一静态分析。余单子 $(D, \varepsilon, \delta)$ 捕获上下文依赖计算：$\varepsilon : D(A) \to A$ 从上下文中提取当前值，$\delta : D(A) \to D(D(A))$ 复制上下文以供嵌套访问。环境余单子 $D(X) = E \times X$ 建模对固定环境 $E$ 的依赖；流余单子 $D(X) = \mathbb{N} \to X$ 建模对时序数据的依赖。

\subsubsection{分次协效应}
\label{sec:2.2.2}

为实现更细粒度的追踪，分次协效应系统使用预序半环 $\mathcal{S} = (S, \leqslant, +, \times, 0, 1)$ 作为协效应代数 \cite{ref33}，这一方法后来由 Gaboardi 等人 \cite{ref19} 与分次效应统一。$S$ 的元素标注每个变量绑定以量化其使用：0 表示未使用，1 表示线性使用，$n$ 表示有界使用，$\infty$ 表示无限制使用。半环运算分别按顺序（$\times$）和并行（$+$）组合协效应，从而在统一的代数框架 \cite{ref37} 内实现精确的资源追踪、灵敏度分析 \cite{ref34} 和信息流控制 \cite{ref35, ref36}。

\subsection{与动态可组合性的关系}
\label{sec:2.3}

效应与协效应系统沿两个互补方向组织关于计算的推理：效应描述计算如何修改其环境，而协效应描述计算如何依赖其环境。这两个方向对应于第 \S\ref{sec:1} 中所识别的动态可组合性的两个维度：
\begin{itemize}
\item 时间可组合性要求组件对共享环境的修改在卸载时可逆。相关的效应是有状态的效应，它们持久地变换该环境；撤销这种变换要求它具有逆。
\item 空间可组合性要求组件间依赖被声明并响应式地管理。此类依赖正是协效应所捕获的内容，管理它们就是将每个依赖与环境的供给进行匹配。
\end{itemize}

然而，经典的效应与协效应系统是静态工具：效应在词法固定的作用域内被追踪，并由编译时处理器处置；协效应标注依据执行前确定的上下文进行验证。相比之下，动态组合要求这些保证对在运行时到达和离开的组件成立，且所面对的上下文持续演化。任何固定的词法作用域都无法界定部署后加载的插件；任何编译时上下文都无法预见运行时配置所产生的依赖。

这促成了视角的转变：我们不是用更多标注来扩展静态类型系统，而是将效应与协效应的概念结构具体化，使得运行时可以直接对其进行操作，从而动态地建立这些系统在静态时所提供的保证。